\documentclass[11pt]{article}
\usepackage[T1]{fontenc}
\usepackage[utf8]{inputenc}
\usepackage{lmodern}
\usepackage{geometry}
\usepackage{amsmath}
\usepackage{amssymb}
\usepackage{amsthm}
\usepackage{booktabs}
\usepackage{hyperref}
\usepackage{listings}
\geometry{margin=1in}

\lstdefinestyle{lean}{
  basicstyle=\ttfamily\small,
  breaklines=true,
  columns=fullflexible,
  keepspaces=true,
  showstringspaces=false,
  frame=single,
  framesep=4pt,
  xleftmargin=2pt,
  aboveskip=8pt,
  belowskip=8pt,
  literate=
    {⟨}{{$\langle$}}1
    {⟩}{{$\rangle$}}1
    {→}{{$\to$}}1
    {←}{{$\leftarrow$}}1
    {↔}{{$\leftrightarrow$}}1
    {¬}{{$\neg$}}1
    {∧}{{$\wedge$}}1
    {∨}{{$\vee$}}1
    {∀}{{$\forall$}}1
    {∃}{{$\exists$}}1
    {∈}{{$\in$}}1
    {∉}{{$\notin$}}1
    {⊆}{{$\subseteq$}}1
    {⊂}{{$\subset$}}1
    {⊕}{{$\oplus$}}1
    {⊥}{{$\perp$}}1
    {∘}{{$\circ$}}1
    {∪}{{$\cup$}}1
    {∩}{{$\cap$}}1
    {≤}{{$\leq$}}1
    {≥}{{$\geq$}}1
    {≠}{{$\neq$}}1
    {×}{{$\times$}}1
    {·}{{$\cdot$}}1
    {∅}{{$\emptyset$}}1
    {∂}{{$\partial$}}1
    {Γ}{{$\Gamma$}}1
    {Δ}{{$\Delta$}}1
    {Ω}{{$\Omega$}}1
    {σ}{{$\sigma$}}1
    {τ}{{$\tau$}}1
    {φ}{{$\varphi$}}1
    {ψ}{{$\psi$}}1
    {α}{{$\alpha$}}1
    {β}{{$\beta$}}1
    {₀}{{$_0$}}1
    {₁}{{$_1$}}1
    {₂}{{$_2$}}1
    {₃}{{$_3$}}1
    {₄}{{$_4$}}1
    {₅}{{$_5$}}1,
}
\lstset{style=lean}

\newtheorem{theorem}{Theorem}
\newtheorem{proposition}{Proposition}
\newtheorem{lemma}{Lemma}
\newtheorem{corollary}{Corollary}
\newtheorem{remark}{Remark}
\theoremstyle{definition}
\newtheorem{definition}{Definition}

\newcommand{\GG}{\mathbb{G}}
\newcommand{\FF}{\mathbb{F}_2}
\newcommand{\Wz}{W_0}
\newcommand{\Cl}{\mathrm{Cl}}
\newcommand{\KC}{\mathrm{KC}}
\newcommand{\LKR}{\mathrm{LKR}_{\mathrm{in}}}
\newcommand{\colorset}{\{r,b,p\}}

\title{A Living Machine-Checked Companion to Goertzel and Fable's\\
Reductive--Compositional Route to the Four-Color Theorem:\\
\large Proven Structure, Refuted Shortcuts, and the Closure Frontier}
\author{The MeTTapedia formalization effort}
\date{Living revision: 13 August 2026}

\begin{document}
\maketitle
\sloppy

\begin{abstract}
This document is intentionally cumulative.  It preserves earlier audits of
Goertzel's v23 proposal, including negative results and abandoned repairs,
while the controlling status is now the current Goertzel/Fable
reductive--compositional programme described in
\emph{reductive\_vs\_covering (2)}, \emph{compositional\_proof\_note},
\emph{4cp\_new\_theory}, and \emph{4cp-lean\_playbook}.  Where that package
supersedes v23, the newer text controls.  Section~\ref{sec:current} records the
current kernel-checked boundary: a large conditional descent tower is present,
but the source-local-to-global realization of the required separated annular
crosscuts, the concrete reductive-system instance, and the verified finite base
remain open.  Accordingly this companion does not claim a new proof of the
Four-Color Theorem.

This is a living formalization companion to Ben Goertzel's manuscript \emph{A
Human-Readable, Spencer--Brown/Kauffman-Style Proof of the Four-Color Theorem}
(v23), which organizes a candidate human-readable proof of the Four-Color Theorem
(4CT) around three interacting pillars: (A) an annular difference-\emph{spanning}
statement for a zero-boundary space $\Wz(H)$ (Theorem~4.9); (B) a CAP5-free
normal form with radius-1 collar locality; and (C) an \emph{executability} layer
that, refusing a false additivity of Kempe switches, reduces between-region Kempe
connectivity to a single \emph{Local Kempe Reachability with fixed input}
($\LKR$) property for one canonical three-cell collar gadget $\tau$ (Lemma~8.14),
composes it along the collar chain (Lemmas~8.15, 8.18, Cor.~8.19), and assembles
the wide annulus collar-by-collar (\S9.2). We record what the MeTTapedia
formalization has machine-checked of this route, with an honest ledger separating
what holds outright, what holds modulo one named remainder, and what is open, and
we foreground the places where the formalization \emph{refutes} a step rather than
discharging it.

Machine-checked outright (no \texttt{sorry}; axioms $\subseteq
\{\texttt{propext},\texttt{Classical.choice},\texttt{Quot.sound}\}$): finite
encoded-model audits of the keystone $\LKR$ property (Lemma~8.14) and two-color
\emph{transparency} (Lemma~8.15), together with completeness of the 192-state
incidence enumeration --- graph-level Kempe-step soundness remains open; a
proper-subspace strictness theorem against a larger unrestricted
boundary-zero chain space, with four uniform peel oracles refuted and
conditional need-based span infrastructure --- v23's actual cycle-space
inclusion (Theorem~4.9) remains open; and a length-2/3
radial-closing-edge fiber census for \S9.2 Step~3 --- the manuscript's
all-outer-stub, arbitrary-length statement remains open. The Gate-2 composition
(Lemmas~8.18/8.19) is \emph{reduced} --- the manuscript's own pointwise Step-1
preparation is refuted and replaced by a frontier-automaton / fibration
architecture --- to one live pin combining representative real-fiber
connectivity audits with DFA-mode sufficiency.

The formalization also produced a decisive negative in the encoded model.
\S9.2 Step~4 --- the manuscript's reinterpretation of collar Kempe words as
words of the ambient annulus --- fails in the encoded two-layer tau-chain
model: a certified, collar-input-avoiding switch whose ambient component
escapes the inner boundary, re-enters, and changes the collar restriction. Its
realization as a simple planar, cyclically-five-edge-connected actual collar is
a named open obligation, so this is a regression countermodel against the
composition \emph{seam} as written, not a graph-level refutation --- and not a
refutation of 4CT; its salvage (a normal-form confinement argument, or a
face-collar redesign) is open. We make no claim about the Four-Color Theorem
--- which is of course a theorem, via the Appel--Haken and
Robertson--Sanders--Seymour--Thomas computer-assisted proofs. The
human-readable v23 route remains \emph{open}, with its foundational layer
audited in an encoded model, and its most dangerous assembly seam pinned down
as broken in that model.

A third development round turned to the v24 working material (\emph{Deserts,
Pivots, and the Necklace}), which replaces v23's \S6--7. The v24 validation
toolchain and its numerical claims were reproduced independently: the $5412$
and $31704$ annular censuses with their Kempe fibers, the fullerene bases,
the pentagon seals, the desert and strand specimens --- including the
$688$-edge strand refuting the package's own Short Strand Conjecture --- and
the GWCO exact base through twelve. On the structural side the development
now carries a rotation-system Tait-coloring layer with exact facial linear
algebra; the all-face discrete Gauss--Bonnet identity
(\texttt{totalFaceCurvature\_eq\_twelve}); a proved \emph{weighted} pentagon
budget correcting the package's at-most-twelve pentagon count; the
locked-pentagon necklace assembly with singleton-strand forcing; a hexagonal
corridor theory (rung classification, clean width-four slabs, oriented
interfaces, Tait transfer, a sharp chirality barrier); congruence-correct
corridor and tube-ring pumping; and a framed-trail layer with
square/digon/rotor fiber calculi. The route's endgame --- fusion chains,
chiral monodromy, Trail Completability with its L1--L10 dependency package,
and a proved Short-Strand replacement --- remains open; the corridor and
pumping infrastructure is the active replacement search.
\end{abstract}

\tableofcontents

\section{Current controlling status: the reductive--compositional programme}
\label{sec:current}

\paragraph{Source precedence and route identity.}
The present development formalizes the reductive route: a sufficiently large
counterexample contains a finite-state corridor, two equal interface states
permit a physical splice, and the splice produces a strictly smaller
counterexample.  This is not the classical unavoidable-set, reducibility, and
discharging architecture.  No configuration catalogue, discharging system,
hubcap family, or per-configuration ring-coloring table is admissible as a
substitute.  The current Fable manuscripts listed in the abstract govern this
section; v23 remains a historical and mathematical source where it is
consistent with them.

\subsection{The idea in plain language}

Imagine walking along a long annular corridor and recording, at every
transversal, only what the rest of the graph can observe through that cut.  A
full profile records boundary colors, tracked-color connectivity, face
continuation, fragment-to-port incidence, and capped face progress.  There are
only finitely many such records.  A sufficiently long walk therefore repeats a
record.  If the two equal records occur on actual separated simple
transversals, the portion between them should be removable: the two exposed
ends carry the same information, so the exterior should not notice the splice.

The word ``actual'' is the geometric heart of the proof.  Finite-state algebra
can prove what follows \emph{from} a suitable pair of cuts; it cannot construct
those cuts merely by naming them.  The current frontier is precisely the
construction of a long, end-capped, source-positioned pair of simple dual
crosscuts from the Cell--3 corridor geometry of the minimal-counterexample
annulus.

\subsection{What is presently machine-checked}

The following items are green Lean developments.  Their scope matters: several
are consumers of a crosscut pair rather than constructors of one.

\begin{center}\small
\begin{tabular}{@{}p{0.34\linewidth}p{0.57\linewidth}@{}}
\toprule
Layer & Kernel-checked content and boundary \\
\midrule
Finite profiles and L7 & A generic finite carrier, decidable equality, exact
cardinality formulae, and finite-profile pigeonhole machinery. \\
Relational transfer & Open-tangle relations, heterogeneous transfer words,
composition, reachable-set pumping, and colorability-reflection consumers.
The primitive transfer is relational; determinism is not assumed. \\
Literal profile semantics & The five profile coordinates are related to actual
Tait colorings.  Tracked connectivity and occurrence-sensitive facial progress
are updated under one regional seam condition. \\
Portal completeness & Both tracked components and face occurrences escaping a
literal prefix are routed to named interface portals. \\
Splice consumers & A literal Tait-coloring gluing equivalence, full-profile
diagonal audits, and strict output-size decrease.  These consume the missing
physical realization and reflection theorem. \\
Source structure & Sector Alternation, two-rail separation, local Cell--3
formation, consecutive side-interface matching and orientation, and both
local rail-edge identifications, with their exact successor cyclic positions,
packaged as a canonical one-step rail cell.  Consecutive surviving side slots
now also give a literal exterior facial-dual rail edge using only cubicity at
the displayed interior corner, and the finite cyclic calculation classifies
every nonadjacent cell as the exact rail shape (0+2), (1+1), or (2+0).
Each case now constructs two literal simple facial-dual walks, each of length
at most two and with combined length exactly two; the two local supports are
disjoint.  Every face in either walk is additionally certified adjacent to
the Cell--3 center that generated it.  Combining this provenance with the
boundary-clean geodesic exclusion now proves all four track combinations
disjoint whenever their centers are at least three positions apart.  The two
bounded neighboring gaps remain a finite local obligation.  For two
consecutive cells, that obligation is now isolated as the proposition
\texttt{CommonNeighborsExact}: the only common full-dual neighbours of the two
centres are the two named flank faces of their shared rung.  Assuming exactly
this proposition, the literal successor rails rebase at the preceding exits,
append to two simple facial-dual paths, and retain mutual support
disjointness.  This is a complete conditional two-cell append theorem, not
yet a derivation of the condition from the source annulus.  A literal
successor theorem carries each outgoing flank to an
incoming flank of the next cell on the same ambient edge and exterior face,
then returns one dependent successor package containing the next outgoing
flanks and the exhaustive local shape certificate.  The literal walks are
generated definitionally from that retained certificate, rather than returned
by an independent existential choice.  Thus a later induction cannot pair a
case certificate with unrelated paths or forget which of the three literal
cases generated them.  Lengthwise simplicity and end-cap
attachment remain open.  These are local formation facts, not yet a global
annular crosscut pair. \\
\bottomrule
\end{tabular}
\end{center}

This layer also contains a useful discipline theorem: equality of boundary
color words alone cannot replace equality of full profiles.  A two-state
machine-checked counterexample admits an off-diagonal transition while having
matching projected color support and no full-profile diagonal.  Connectivity
and face data are therefore mathematical content, not bookkeeping.

\subsection{The live geometric gate}

The type \texttt{SeparatedAlignedSimpleDualCrosscuts} already states the right
output: two aligned simple dual crosscuts, disjoint transverse supports, and a
nondegeneracy condition.  Much of the splice development is parametrized by
such an object.  What is not yet in the tree is a theorem producing the required
long/end-capped instance from the actual Cell--3 annular carrier and its
selected source geometry.

\paragraph{Source-alignment correction (13 August 2026).}
Earlier local-rail modules used ``L1'' as a convenient internal tag.  That
was misleading.  In the current Fable manuscript and playbook, L1 is the
bulk statement that a sufficiently large frontier instance contains a long
hexagonal ladder corridor.  The bounded common-neighbour premise used to
append two locally constructed rails is not stated as Fable's L1, and no
current-source theorem yet identifies it as the construction of L1's
corridor.  The original v23 \emph{boundary locality} argument is historical
evidence about radius--one collars, but it assumes confinement and
bounded-interaction claims that the corrected Fable work records as refuted.
It cannot license the present rail assembly.

Accordingly, the local rail developments below are retained as elementary,
kernel-checked geometry and as a conditional assembly interface; they do not
count as a source realization or as progress through a Fable flag.  The next
source-faithful question is narrower: which concrete current-source formation
constructs the simple transversals/layer boundaries used by the Cell--3 splice,
and does it entail the rail premise at all?  Until that formation is identified
and proved, neither the L8 radial witness nor the closed-map opening transport
is silently treated as the missing crosscut construction.

There is, however, now a direct handoff for the current Lean clean-block
extraction which narrows this question further.
A \texttt{ClosedWebAtGoodWord.Instance} already contains the actual
\texttt{ClosedWebAnnularEmbedding} and its \texttt{AnnularFrontierGeometry}.
The theorem
\texttt{exists\_sourceRealizedBoundaryCleanOrbitHexCorridor\_of\_weightedCleanBlock}
therefore turns the manuscript's explicit weighted size premise into a
boundary-clean corridor together with the literal source geodesic and selected
block.  It does not reopen the source through a cap-deletion model.  This is a
real conclusion of the current Lean extraction, conditional only on the stated
quantitative bound.  It is not yet Fable's full L1 ladder corridor: the latter
also needs the local rung-incidence classification used to produce simple
transversals.  One initially apparent part of that gap has now been removed
without imposing a closed-minimal hypothesis.  The new
\texttt{SelectedCorridorRungs} object chooses one genuine shared primal edge
for each consecutive dual step, rather than requiring all pairs of faces to
share at most one edge.  Lean then proves that the selected incoming and
outgoing rungs of every internal Cell--3 hexagon occupy distinct,
non-adjacent positions on its six-dart boundary.  The proof uses only the
source carrier's interior-face simplicity, local cubicity, and geodesic
separation.  This is a local rung-placement result, not a proof of the
remaining finite adjacency classification, a long separated pair of
crosscuts, or a splice.

Three nearby results delimit the gap.

\begin{itemize}
  \item Generic dual connectedness constructs one hole-to-hole crosscut.  It
        does not produce a separated pair tied to the selected corridor.
  \item The local Cell--3 layer produces a separated local interface package,
        but its endpoints are not the required distant annular endpoints.
  \item Consecutive six-edge aligned slab-boundary walks are proved
        non-disjoint: they share their intervening dual interface.  They are
        useful local windows, not the two radial cuts of the splice.
  \item The source-facing \texttt{LongRadialSectorWitness} already supplies
        two vertex-disjoint \emph{primal} radial paths, distinct anchors on
        both boundary components, sector alternation, and a quantitative
        length witness.  It does not itself supply facial-dual crosscuts, and
        no module presently relates this witness type to
        \texttt{SeparatedAlignedSimpleDualCrosscuts}.
  \item The coverage clause is now explicit in the ambient graph:
        \texttt{interior\_vertex\_mem\_ambientRadialPath\_union} transports
        \texttt{RadialPathPair.cover\_support} to prove that every interior
        vertex lies on one of those two radial paths.  This is precisely the
        ``jointly covering every internal vertex'' part of Addendum~XXII; it
        still does not construct the two complementary sector sides.  It does
        mean that any later sector side proved disjoint from both paths is
        immediately free of interior vertices.
\end{itemize}

The last item is a scoped refutation of a tempting Lean construction, not a
refutation of Goertzel's proof.  It says that the separated cuts must be
assembled at the correct global scale rather than identified with adjacent
local windows.  Nor is the live problem a general Menger-type existence
theorem.

There is an important source-level separation here.  In Addendum XXVII and
the formalization playbook, the radial-path witness is the L8 input to the L6
length/depth dichotomy, whereas the boundary-clean hex corridor is the
independent L1 output.  The manuscripts do not instruct us to dualize the L8
witness into the L1 corridor crosscuts.  The absence of a Lean module relating
the two types is therefore not by itself a missing source theorem, and we do
not silently identify them.

The literal Cell--3 rail construction is a Lean-side exploratory candidate
for a physical corridor realization; the current sources do not yet identify
it with Fable's L1 construction.
Its remote noncollision is proved, and its adjacent append has been reduced to
\texttt{CommonNeighborsExact}.  Lean now proves that this face-level statement
is equivalent to an edge-level statement: for every common neighbour, the
canonical primal edge from the left centre to that face is one of the two
literal flank edges of the shared rung.  The proof uses the fact that an edge
of a planar cellular map is incident with at most two face orbits, together
with a new exact identification of each literal flank edge with the canonical
shared edge of its centre and side face.  Thus the bounded local gap is not a
second, independent classification of faces; it is the transport of one
closed-map primal-edge classification into the literal open Cell--3 carrier.
This equivalence is a reduction of the obligation, not its discharge.

An extra common neighbour would form precisely
the small separating/nonfacial dual triangle excluded by v23's
historical minimal-counterexample radius--one collar normal form.  The present open-annular
\texttt{Instance}, however, records well-formed boundary stubs, coloring,
closure, and the good word, but not the closed minimal-counterexample
provenance from which that bounded-interaction fact is derived.  That old
normal form is not a current-source justification for adding this
classification to the Cell--3 instance.  Any future use needs a separate,
current-source formation theorem; until then the result is only a conditional
Lean interface.

The closed half of that argument is now machine-checked.  Given three
pairwise adjacent faces in a graph-backed vertex-minimal Tait counterexample,


\begin{center}
\texttt{exists\_singleton\_primalCut\_of\_dual\_triangle}
\end{center}


constructs their literal length-three facial-dual cycle and proves that its
three crossed primal edges separate a component containing exactly one
vertex; moreover, the crossing-edge set of that singleton component is
exactly the three edges read from the dual cycle.  Lean now also extracts the
local consequence used by the radius--one collar description:

\begin{center}
\texttt{exists\_vertex\_mem\_all\_crossed\_edges\_of\_dual\_triangle}.
\end{center}

The consumer-facing theorem

\begin{center}
\texttt{exists\_vertex\_mem\_three\_shared\_edges\_of\_dual\_triangle}
\end{center}

then identifies the three canonical pairwise shared primal edges themselves
and proves that one ambient vertex belongs to all three.  Thus the dual
triangle is the link of one primal vertex, rather than an arbitrary separating
triangle.  This is the closed-map bounded-interaction classification needed by
the source.

The closed classification is now complete in
\texttt{GoertzelV24MinimalCorridorCommonNeighbors}.  Its theorem

\begin{center}
\texttt{commonNeighborEdge\_eq\_outgoing\_flank}
\end{center}

starts with a literal hex-corridor placement in the closed graph-backed
minimal counterexample.  For any face adjacent to both a corridor face and
its successor, the theorem proves that the canonical shared primal edge from
the first centre is exactly one of the two face-cycle edges immediately
flanking the outgoing rung.  The proof combines the common vertex extracted
from the dual triangle with cubicity and two-sided face incidence: at either
endpoint of the outgoing rung, the face has exactly two incident boundary
edges, namely the rung and its immediate flank.  Since the target edge is not
the rung, it is the flank.  Thus the bounded local classification is no longer
an assumption on the closed map.

What remains is transport, not another closed classification.  The three
faces, the outgoing rung, and its two flank edges must be carried through the
actual source opening and identified with the literal Cell--3 objects.  Only
then does this theorem discharge \texttt{CommonNeighborsExact} on the open
carrier.  The latter has degree-one boundary stubs and is provably incompatible
with the closed cubic minimal-counterexample structure, so adding minimality
as a field of the open instance would be inconsistent, not a shortcut.  The
remaining statement must construct the explicit closed-map-to-open-annulus
formation and preserve the relevant interior faces and shared boundary edges.

The strictly local open-carrier implication is now machine-checked as well.
The predicate \texttt{CommonNeighborVertexIncidence} asks only that, after the
opening, the outgoing rung and the canonical edge from the centre to each
common neighbour still share one primal endpoint.  The theorem

\begin{center}
\texttt{commonNeighborEdgesExact\_of\_commonNeighborVertexIncidence}
\end{center}

uses cubicity at that retained endpoint, edge-simplicity and local
two-sidedness of the displayed interior face, and parity of face-boundary
incidence to force the second edge to be one of the two immediate flanks.  A
second theorem then derives \texttt{CommonNeighborsExact}.  Thus the open
classification no longer needs the entire closed minimality package: it needs
one concrete incidence fact preserved by the source opening.  This is a
reduction of the formation gap, not its solution.  The incidence predicate is
not assumed by the source instance and the new module does not construct the
Cell--3 opening that must establish it.

The generic opening calculation needed by that construction is now available
in \texttt{GoertzelV24OpenRegionEndpointIncidence}.  If two ambient edges meet
at a vertex and both edge occurrences lie on fully retained faces, their
computed literal-open edges meet at the corresponding retained vertex.  The
vertex-retention fact is derived from full face retention rather than supplied
as an extra premise, and both open edges live in the same fixed opened rotation
system.  Consequently the remaining incidence transport is no longer an
abstract claim that opening ``preserves local geometry'': it is the concrete
source formation, the full-retention proof for the displayed Cell--3 faces,
and their already named edge identifications.

The closed classification and the generic opening calculation have now also
been joined in one constructive theorem.  In
\texttt{GoertzelV24OpenRegionMinimalDualTriangleIncidence}, the theorem

\begin{center}
\texttt{exists\_common\_open\_endpoint\_of\_retained\_dual\_triangle}
\end{center}

starts from a facial-dual triangle in the closed graph-backed
vertex-minimal counterexample and assumes that its first ambient face is
wholly retained by a literal opening.  It does not accept two arbitrary open
edge coordinates.  Instead, it extracts the two canonical closed shared
edges, constructs their actual positions on the first face cycle, computes
the corresponding open edges, and returns a retained open vertex incident
with both.  Hence the generic cap-deletion question at this corner is settled:
under full-face retention, the common primal corner survives as literal open
incidence.  The remaining burden is source-specific and deliberately visible:
construct the Cell--3 opening from the closed map, prove the selected L1 face
is fully retained, and identify its named rung/flank coordinates with these
constructed positions.  The theorem neither assumes nor constructs that
formation, and it does not by itself prove
\texttt{CommonNeighborsExact} or the separated crosscut pair.

The full-retention premise in that theorem is no longer a purely qualitative
condition.  The generic cubic-map lemma in
\texttt{GoertzelV24CubicFaceVertexSeparation} proves that two distinct
facial cycles meeting at any primal vertex must be adjacent in the full
facial dual.  Therefore two distinct dual-nonadjacent faces have disjoint
primal vertex supports.  Taking the literal opening to retain precisely the
complement of the vertices visited by a chosen cap face gives the constructive
corollary

\begin{center}
\texttt{faceFullyRetained\_compl\_orbitFaceVertices\_of\_not\_adj}.
\end{center}

Thus any corridor face proved boundary-clean---distinct and dual-nonadjacent
from the pentagonal cap---is automatically wholly retained by cap deletion;
no separate vertex-by-vertex retention assumption is needed.  This settles
the generic closed-map side of that premise.  It does not select the source
pentagon, prove that the displayed Cell--3 faces are nonadjacent to it, build
the two-hole annulus, or realize the long separated crosscuts.  Those are the
remaining source-specific L1 formation statements.

The next bounded L1 collision has now been classified on the closed map as
well.  Geodesicity rules out any dual edge between corridor faces two steps
apart, but by itself it does not say that their displayed middle face is their
only common neighbour.  If an additional common neighbour exists, the two
length--two routes form a simple dual four-cycle.  In
\texttt{GoertzelV24MinimalDualFourCycleClassification}, cyclic five-edge
connectivity converts this cycle into the already established adjacent-pair
primal collar and supplies a chord joining opposite cycle positions.  The
endpoint chord is impossible by geodesicity, so Lean proves

\begin{center}
\texttt{middle\_adj\_other\_of\_two\_common\_neighbors}.
\end{center}

Its corridor specialization says that every such extra common neighbour must
also be adjacent to the displayed middle face.  Thus the residual two-step
collision is not an arbitrary global interaction: it is forced into the
finite adjacent-pair collar geometry named by L1.  This is a genuine
classification theorem, not a conditional rail-assembly wrapper.  It does
not yet prove that the extra face is impossible, discharge
\texttt{CommonNeighborsExact} in the cut-open Cell--3 carrier, or construct
the source opening; those are the next finite local and formation steps.

The theorem is already consumed by the closed corridor API rather than left
as an isolated candidate.  The derived result
\texttt{twoStepCommonNeighbor\_eq\_middle\_or\_outgoing\_flanks} combines it
with the two adjacent dual-triangle classifications: a common neighbour at
index distance two is either the middle corridor face itself, or its shared
edge from each adjacent centre belongs to that centre's named flank pair.
Consequently all rail-collision distances now have a distinct proved status:
remote collisions are impossible by geodesicity, consecutive collisions are
classified by one flank pair, and distance-two collisions are trapped between
the two adjacent flank pairs.  The remaining closed local task is the finite
compatibility of those two flank classifications, not an unbounded path
existence theorem.

The corresponding generic opening step is now proved in
\texttt{GoertzelV24OpenRegionMinimalDualFourCycleClassification}.  Its theorem

\begin{center}
\texttt{open\_middle\_adj\_open\_other\_of\_retained\_two\_common\_neighbors}
\end{center}

starts with the closed graph-backed minimal-counterexample four-cycle and a
literal vertex-region opening.  If the middle and extra facial orbits are
fully retained, the forced middle--extra dual edge remains an edge of the
opened facial dual.  The proof uses the exact retained-face adjacency
equivalence, so it cannot manufacture an adjacency at a boundary stub or lose
one by comparing only cardinalities.  This is the four-cycle analogue of the
retained dual-triangle corner theorem.  It is still a generic transport
theorem: the source-specific Cell--3 formation must identify its literal faces
with these open images and establish their full retention.

One tempting strengthening is false and is excluded from the assembly design.
A middle Cell may have a zero-length rail on one track (the
\texttt{forwardTwo} or \texttt{forwardFour} shape), so an incoming flank can
equal the same-track outgoing flank.  Such a repetition is the legitimate
connector removed by the tail convention in walk append, not a self-
intersection.  The remaining finite theorem must therefore show that every
distance-two overlap is this same-track connector; it must not claim that all
distance-two overlaps are impossible.

This carrier distinction also rules out a superficially convenient adapter.
The rotation system in a well-formed \texttt{SourceTrail.AnnularEmbedding} is
already the cut-open framed graph and is proved not globally cubic, whereas
the minimal-counterexample theorem concerns the original closed cubic map.
One may not assume both structures on that same carrier and call the resulting
conditional theorem a formation result.  This does not leave the Cell--3 L1
corridor unformed: its source instance already stores a
\texttt{ClosedWebAnnularEmbedding} and the required annular geometry, and the
weighted clean-block theorem now returns a source-realized dual geodesic from
those fields.  The missing bridge includes the source's ladder/rung-incidence
classification and then the construction of long separated transverse
crosscuts needed by the splice.  The generic
cap-deletion construction remains relevant only if a later source theorem
identifies that particular opening with the Cell--3 carrier; it is not silently
inserted as the current L1 construction.

The first generic piece of this formation is now constructed in
\texttt{GoertzelV24OpenRegionFaceTransport}.  Let a closed ambient face orbit
be wholly based at vertices retained by the literal open-region construction.
Lean proves that following the ambient face permutation never enters a fresh
boundary stub, that the old-dart embedding commutes with every face step, and
that the complete ambient and open dart cycles are equivalent.  All transported
faces now live in one fixed literal-open rotation system, rather than in
separately rooted copies.  In that common map Lean proves that distinct retained
ambient faces remain distinct, that equal ambient edge occurrences compute the
same open edge even when they belong to different faces, and hence that a
shared ambient boundary edge remains a genuine facial-dual adjacency after
opening.  Conversely, every open dual adjacency between two such transported
faces is pulled back to a shared ambient edge occurrence.  Consequently Lean
now proves the graph-level equivalence

\begin{center}
\texttt{openFaceOrbit\_adj\_iff\_ambientFaceOrbit\_adj},
\end{center}

so the literal opening induces exactly the same adjacency relation on this
family of retained face orbits: it neither loses an ambient dual edge nor
creates a spurious one.  Thus an untouched face is transported as an actual finite orbit
without splitting or identifying its old boundary edges, not merely identified
by a matching cardinality or assumed through a wrapper.  This is generic
carrier plumbing for any future cap-deletion formation.  It neither supersedes
the direct Cell--3 L1 handoff nor constructs the separated facial-dual
crosscut pair from its source-realized corridor.

The next formation step is now explicit in
\texttt{GoertzelV24OpenRegionHexCorridorTransport}.  First, Lean proves that a
fully retained two-sided face has exactly the same boundary-support cardinality
after opening.  This is not inferred from two matching numbers: the retained
face-cycle equivalence is combined with an edge-level injectivity proof that
reflects equality of opened edges back to the ambient face.  It then constructs
\texttt{openOrbitHexCorridorSkeleton}: any closed hex corridor all of whose
selected faces are fully retained becomes an actual hex corridor in the one
computed open rotation system, with the same length, injective face indexing,
hexagonal boundaries, consecutive adjacencies, and separated non-adjacencies.
Thus the generic transport of an L1 corridor is complete.  It remains
conditional on a cap/collar opening being the relevant source carrier; the
current Cell--3 clean-block path already has its own direct annular
realization.  Neither route supplies the source's full ladder incidence or
the long, end-capped crosscut pair, which remain the live geometric
obligations.

There is also a cap-specific constructor,
\texttt{openOrbitHexCorridorSkeleton\_compl\_cap}.  In a cubic, cyclic,
two-sided closed map, suppose every selected corridor face is distinct and
dual-nonadjacent from a chosen cap face.  The previously proved cubic
face-separation theorem then shows that deleting precisely the cap-face
vertices fully retains every corridor face, and the generic transport
constructor produces the open corridor.  This reduces the cap-opening part of
L1 to a source-visible premise---the displayed dual non-adjacencies---rather
than a vague claim that opening should preserve the picture.  It still does
not select the source cap or prove those non-adjacencies.

Retaining only the selected corridor faces is not sufficient for the clean
L1 interface.  A neighbouring open face could otherwise be a cut fragment
whose facial cycle was closed through a fresh boundary stub.  The module
\texttt{GoertzelV24OpenRegionCleanHexCorridorTransport} closes this loophole.
Its one-sided neighbour-pullback theorem starts from an actual edge shared by
a retained face and an arbitrary open neighbour, recovers the opposite
ambient dart, and uses the two-face incidence bound to prove that the open
neighbour is exactly the image of that ambient face.  Consequently
\texttt{openCleanOrbitHexCorridorSkeleton} transports an entire clean
corridor whenever its radius-one ambient neighbourhood is fully retained.
The cap specialization packages the geometric premise as
\texttt{CorridorOneRingAvoidsFace}: every centre and every adjacent face is
distinct and dual-nonadjacent from the cap.  Cubic face separation then
constructs all retention proofs automatically.  Thus boundary cleanliness is
preserved as a theorem about actual face orbits, rather than inferred from a
picture or from center-face retention alone.

The adjacent common-neighbour classification has now been upgraded from its
edge form to its literal face form.  The same closed module defines the face
across a side-slot dart using only the ambient rotation system and proves

\begin{center}
\texttt{commonNeighbor\_eq\_outgoing\_sideFace}.
\end{center}

Thus a face adjacent to both consecutive corridor centres is not merely known
to use one of two flank edges: two-sided facial incidence proves that it is
exactly the opposite-dart face at one of those two flank positions.  This
generic definition is deliberately independent of the later closed-web
namespace; the proof uses the closed minimal-counterexample hypotheses only
where they supply the edge classification.

The new module
\texttt{GoertzelV24OpenRegionMinimalCorridorCommonNeighbors} then joins that
closed theorem to the boundary-clean opening.  An arbitrary open common
neighbour is pulled back, from each of the two corridor centres, to a fully
retained ambient neighbour.  Injectivity of retained face-orbit transport
shows that the two pullbacks represent the same ambient face, after which the
closed theorem classifies it as one of the two literal side faces.  The result
\texttt{exists\_ambient\_side\_rep\_of\_open\_commonNeighbor} therefore proves
the face-level bounded interaction through the computed opening without
assuming cubic minimality of the noncubic open carrier, invoking a general
disjoint-path theorem, or matching carriers by cardinality.  It still leaves
the source specialization visible: the actual Cell--3 cap/corridor selection
must establish radius-one retention and identify the generic transported
side faces with its named local rail data.

The representative choice in that statement has now also been removed.
Fully retained roots on one ambient facial orbit have equal images in the
fixed opening, by

\begin{center}
\texttt{openFaceOrbit\_eq\_of\_dartOrbitFace\_eq}.
\end{center}

Consequently \texttt{open\_commonNeighbor\_eq\_sideFaceImage} states the
consumer-facing result directly: every open common neighbour is one of the
two explicit open facial orbits rooted at the opposite darts of the outgoing
flanks.  This is stronger than retaining an unspecified ambient representative
and then identifying its quotient face.  It remains a generic literal-opening
theorem, however; it does not assert that the computed opening and the
source-local closed-web annulus are definitionally the same carrier.  The next
source formation adapter must prove that identification or an isomorphism
preserving these named roots.

The opened rotation is now also backed by an actual computed simple graph in
\texttt{GoertzelV24OpenRegionGraphBacking}.  Its primal adjacency is defined
from the opened dart endpoints.  Lean proves that mapping an opened dart to
its ordered pair of endpoints is bijective: internal darts are recovered from
the two retained endpoints, while a fresh boundary vertex remembers the exact
exposed ambient half-dart.  Transporting the opened vertex rotation across
this equivalence produces a \texttt{SimpleGraphDartRotation.Data} package on
the computed graph.  The equivalence is proved to commute with edge reversal,
vertex rotation, and hence the facial step.  Thus later Cell--3 formation may
move between the literal-open rotation carrier and an ordinary graph-backed
carrier without silently changing face or edge semantics.  Moreover, if the
ambient rotation is genuinely cyclic at every vertex, the opened rotation and
its graph-backed presentation inherit that property: retained vertices keep
their complete cyclic dart fibre, while every fresh boundary vertex has a
singleton fibre.  The same module now constructs the transparent graph
homomorphism from the ambient graph induced on the retained vertices into the
opened graph.  Every fresh boundary vertex is proved adjacent to the retained
endpoint of the exposed ambient half-dart that created it.  It follows that a
connected retained induced region remains connected after opening: the new
vertices are leaves, not unaccounted-for components.  These results discharge
the cyclic-rotation and connectedness fields of a future annular cellulation,
once source-specific formation supplies the corresponding ambient cyclicity
and retained-region connectedness proofs.  This remains generic carrier
plumbing: it neither selects the Cell--3 retained region nor constructs its two
hole faces, planar cellulation, L1 corridor, or separated crosscuts.

There is a second, deliberately separate formation result in the tree.  The
first carrier-changing part of the \emph{single-edge deletion} model is now
constructed, rather than postulated.  This is the formation geometry used by
the missing-edge framed trail of Addendum~XXVIII (Cell--4); it is not the
annular Cell--3 opening of Addendum~XXVII.  The generic permutation lemma

\begin{center}
\texttt{eraseTwoPointsSubtype}
\end{center}

removes the two oriented darts of the distinguished edge from the ambient
vertex rotations: each removed dart becomes fixed in the auxiliary
permutation, and its old predecessor is spliced directly to its old successor.
Restricting to the complement and transporting along the exact dart
equivalence gives

\begin{center}
\texttt{deletedEdgeData} :
\texttt{Data G} $\longrightarrow$
\texttt{Data (DeletedEdgeGraph G u v)}.
\end{center}

Every surviving dart remains based at the same vertex.  More importantly for
face geometry, Lean proves

\begin{center}
\texttt{deletedEdgeData\_phi\_ambient\_eq\_of\_ne}.
\end{center}

If an ambient face step does not land on either orientation of the removed
edge, the face step computed in the deleted graph maps back to exactly the old
ambient face step.  Thus deletion changes face tracing only where the missing
edge was actually encountered; remote face steps are not reconstructed by an
unproved picture.  This is still only a bare rotation-system construction and
a local transport theorem.  It does not transport planar-geometry fields,
choose the source's missing edge, or finish Cell--4 formation.  More
importantly for the present closure frontier, it does not construct Cell--3's
cap-deleted annulus, container cycles, two holes, or long separated crosscut
pair.  Specializing it to Cell--3 merely because both constructions produce
open graphs would confuse two different source objects.

The Cell--4 deletion result has now been lifted to complete untouched faces in
\texttt{GoertzelV24DeletedEdgeFaceTransport}.  The predicate
\texttt{FaceAvoidsDeletedEdge} says that every dart in one ambient face orbit
avoids both orientations of the missing edge.  Under exactly that condition,
\texttt{deletedFaceCycleEquiv} gives a bijection between the whole ambient
dart cycle and the whole computed face cycle of \texttt{deletedEdgeData}.
Surjectivity is derived from the commuting face step and finite cycle
cardinality; it is not assumed.  This proves exact orbit and face-length
preservation for untouched faces.  The same module now proves that two distinct
untouched ambient faces remain distinct, transports equality of their shared
primal edge in both directions, and packages the result as

\begin{center}
\texttt{deletedFaceOrbit\_adj\_iff\_ambientFaceOrbit\_adj}.
\end{center}

Thus single-edge deletion induces exactly the old facial-dual adjacency relation on the
family of untouched face orbits: it neither loses a shared edge nor invents
one between their images.  This is enough to transport a closed-map dual
triangle through the Cell--4 missing-edge opening once its three faces are
proved untouched.  It is not enough to transport the triangle into Cell--3,
whose annular carrier comes from a different region/cap opening.  The live
Cell--3 bridge must instead specialize the generic open-region transport to
the actual closed-map annulus and prove that every open common neighbour under
consideration is the image of an eligible ambient face.

\paragraph{Selected local rungs.}  The direct Cell--3 L1 handoff no longer
assumes that every pair of interior faces shares at most one primal edge.
That global uniqueness premise belongs to a closed minimal-counterexample
carrier and is incompatible with silently reusing it for the source's open
annular tangle.  The new module
\texttt{GoertzelV24ClosedWebSelectedRungPlacement} instead selects one actual
shared edge for every consecutive pair of faces in a boundary-clean source
corridor.  At each internal hexagon it proves, from local face incidence and
the displayed Cell--3 separation, that the selected incoming and outgoing
rungs occupy distinct positions and cannot be adjacent.  Consequently their
cyclic type is one of the two nonadjacent cases,
\texttt{oneEdgeBetween} or \texttt{opposite}; deleting the two selected
rung positions leaves four side slots and exactly two directed cyclic side
steps.  Each selected slot now determines its literal opposite-dart face,
which Lean proves is facial-dual adjacent to the corridor centre and an
annular interior face.  This is a genuine finite local fragment of flag L1,
and it removes a false global hypothesis from the local construction.  It
does not establish that distinct slots give distinct side faces, construct
simple rails, or produce the required long separated crosscut pair.

The first remaining collision has now been reduced without pretending it is
impossible.  If two distinct selected side slots lead to the same opposite
face, Lean exhibits two distinct interior primal edges shared by that face
and the central hexagon.  Thus the open-carrier question is a finite
two-shared-edge local case, precisely the kind of case Fable assigns to L1;
it is not a license to import the closed minimal-counterexample theorem that
would rule the collision out globally.

The next finite case is now explicit as well.  For two consecutive surviving
side slots around a selected Cell--3 hexagon,
\texttt{selectedPlacementSideFaces\_eq\_or\_adjacent\_of\_forwardStep}
proves exactly one of two alternatives: the opposite-dart faces coincide, or
they are joined by the literal exterior facial-dual edge at their common
locally cubic corner.  The first alternative is deliberately retained and is
fed to the two-shared-edge collision theorem above; the second is the actual
one-step rail edge.  This is the intended finite local classification, carried
out on the source annulus, rather than an unproved assertion that every local
rail step is automatically simple.

\paragraph{Do not confuse L8 radial paths with L1 crosscuts.}  Lean already
proves \texttt{exists\_longRadialSectorWitness\_at\_five}: under its
closed-web hypotheses it supplies two vertex-disjoint \emph{primal} radial
paths, distinct anchors on both boundary components, sector-alternation data,
and a quantitative length alternative.  This is the L8-to-L6 package.  The
splice instead requires \texttt{SeparatedAlignedSimpleDualCrosscuts}, a pair
of simple \emph{facial-dual} transversals satisfying a dual-walk disjointness
condition and the crosscut alignment hypotheses.  No theorem currently
relates these types.  Thus the radial witness is a proved neighbouring asset,
not yet a constructor for the L1 pair; treating it as one by renaming or
duality intuition would hide precisely the missing primal--dual construction.

\subsection{The closure chain}

The remaining dependency order is:

\begin{enumerate}
  \item construct the long, end-capped separated crosscut pair from the actual
        annular Cell--3 corridor;
  \item use it to discharge the literal source
        \texttt{SemanticProfileBridge} or the weaker sufficient
        \texttt{ColorabilityReflection};
  \item instantiate the carrier-changing \texttt{ReductiveSystem}.  Its size,
        reduction, and strict-decrease ingredients exist; its concrete
        counterexample-reflection field does not;
  \item derive an explicit threshold and implement \texttt{BaseVerified} from
        the route's own finite base.  No such instantiation exists at this
        revision;
  \item apply the abstract induction theorem to expose the concrete
        source-facing Four-Color statement.
\end{enumerate}

The v23 three-cell micro-enumeration is base material, not a serial corridor
generator.  Conversely, a raw bound on the entire over-encoded profile carrier
is a finiteness witness, not a practical base census.  The eventual base must
bound the reachable state space and replay compact certificates generated by
the manuscript's own definitions.  If the argument were instead found to
require an externally searched unavoidable catalogue, that would be a route
failure requiring an explicit decision, not an implementation detail.

\subsection{What the programme may explain---and what is not yet a theorem}

The finite-state route suggests an information-theoretic explanation for why a
short informal proof may be unavailable: the reduction can require preserving
a large amount of interface information, and the terminal finite verification
may contain irreducible case information.  This does \emph{not} imply that no
human-friendly proof of 4CT can exist in every proof system.

A formal non-compressibility result must first specify a presentation or proof
class, an encoding, a description-length or verification-complexity measure,
and uniform quantifiers.  Only then can a lower bound be meaningful.  The
future formal endpoint is therefore the strongest such theorem supported by
the compositional transfer presentation, together with an explicit boundary
between proved mathematics and philosophical interpretation.  A large state
count by itself is evidence about one representation, not a universal
impossibility theorem.

\subsection{How to read the rest of this companion}

The remaining sections preserve earlier investigation strata: the v23
three-pillar audit, machine-checked obstruction and repair attempts, the
composition ladder, the transition to the v24/Fable work, and the graveyard of
superseded routes.  They are retained because negative knowledge is expensive
and because deleting it makes old mistakes attractive again.  When an older
status statement conflicts with this section or with the current Fable
manuscripts, this section and those primary sources control.

\section{The problem and Ben Goertzel's three-pillar route}
\label{sec:problem}

\subsection{Tait's reduction and the minimal-counterexample target}

The Four-Color Theorem asserts that every plane triangulation $G$ admits a
proper vertex 4-coloring. Passing to the cubic dual $K=G^{*}$, Tait's classical
equivalence (v23, Thm.~2.3) restates 4CT as: every bridgeless planar cubic
graph admits a \emph{proper 3-edge-coloring}, i.e.\ a map
$\chi:E(K)\to\colorset$ assigning distinct colors to the three edges at each
vertex. The manuscript encodes the three colors as the three nonzero elements
of $\GG=\FF^{2}$, with $r=(1,0)$, $b=(0,1)$, $p=(1,1)$, so that ``properness at
a vertex'' is exactly $\chi(e_1)+\chi(e_2)+\chi(e_3)=r+b+p=0$ in $\GG$, and
``integrating'' an edge-coloring across the planar dual becomes a path-sum in
$\FF^{2}$.

Following Spencer--Brown and Kauffman, the manuscript reformulates colorings as
plane \emph{formations} (red/blue Jordan curves with purple overlaps), defines
Kempe switches as the standard two-color interchange, and proves the
\emph{Parity Lemma} (planar invariance of the mod-2 count of two-color
circuits). The role of \S3 is to reduce 4CT, via Kauffman's parity/primality
argument (Theorem~3.1), to the nonexistence of a minimal counterexample
containing a \emph{between-region annulus} $H$ whose boundary-compatible
colorings are \emph{not} connected by Kempe moves supported inside the annulus.
Thus the entire analytic burden is shifted to a \emph{between-region Kempe
connectivity} statement for an annulus $H\subseteq K$ with outer and inner
boundary cycles $\Gamma_\infty,\Gamma_0$.

\subsection{The three pillars}

The manuscript's conceptual overview (v23, \S1.1) states the argument as three
interacting pillars.

\paragraph{Pillar~A: annular difference spanning (the linear ``accounting''
layer).} Working in a between-region annulus $H$, one encodes boundary
colorings and their differences as vectors in an $\FF$-space. Let $W(H)$ be the
coordinatewise cycle space of $\GG$-chains on $E(H)$ satisfying the Kirchhoff/\allowbreak
even-degree constraints at interior vertices, and let
\[
  \Wz(H) \;:=\; \{\,x\in W(H) : x(e)=0 \text{ for all boundary edges } e\in
  B_1\cup B_2\,\}
\]
be the \emph{zero-boundary} subspace. For a fixed proper coloring $C_0$ of $H$
with Kempe closure $\Cl(C_0)$, the manuscript defines a family $\KC(C_0)$ of
\emph{Kempe-closure (polarized face) generators} $X_f^{ab}(C)$ supported in
depth-1 collars (v23, Def.~4.8), and proves:
\begin{theorem}[Spanning theorem for $\Wz(H)$, v23 Theorem~4.9]
\label{thm:49}
Let $H$ be an annular subgraph of a bridgeless planar cubic graph, with
boundary cycles $\Gamma_0,\Gamma_\infty$. Fix any proper 3-edge-coloring $C_0$
of $H$ and let $\mathcal{G}:=\KC(C_0)$. Then
$\Wz(H)\subseteq\mathrm{span}_{\FF}(\mathcal{G})$; equivalently, if
$z\in\Wz(H)$ satisfies $z\cdot g=0$ for every $g\in\mathcal{G}$, then $z=0$.
\end{theorem}
The manuscript proves Theorem~\ref{thm:49} in strong-dual annihilator form: it
shows $U^{\perp}\cap\Wz(H)=\{0\}$ for $U=\mathrm{span}_{\FF}(\mathcal{G})$, by a
planar orthogonality peeling along a dual spanning forest, after a
``purification'' step (Lemma~4.4) that removes the polarization of the face
generators.

\paragraph{Pillar~B: CAP5-free geometry and radius-1 locality (the geometric
locality layer).} Assuming a minimal counterexample, the manuscript eliminates
the CAP5 separating obstruction and works in a CAP5-free normal form
(Corollary~6.9). A tube-forcing argument then reduces every ``escape'' of an
anchored Kempe chain to a finite three-cell gate enumeration, yielding radius-1
confinement of anchored Kempe cycles together with a cap-aware
disjointness/commutativity package (Lemmas~7.8, 7.9, 7.10). This is what
licenses peeling the wide annulus into radius-1 collars.

\paragraph{Pillar~C: executability without false additivity (the
algebra-to-geometry bridge).} The decomposition $\Delta=\sum v_{f,g}$ provided
by Pillar~A is a statement in $\Wz(H)$; it does \emph{not} by itself imply that
one may \emph{execute} the corresponding mixed color-pair Kempe switches as if
they formed an abelian sum. The manuscript is explicit that cross-pair
interference is real: an $(a,b)$-switch can destroy the $(b,c)$-structure a
later switch needs (v23, \S8.1), and Kempe switches using different color pairs
on shared support need not commute (v23, \S8.3). To execute the decomposition
honestly, the manuscript moves to an operator model: proper colorings form an
orthonormal basis of a Penrose--Kauffman state space, and each Kempe switch is
an involutive linear operator (generally noncommuting with switches of other
color pairs). Using radius-1 locality, between-region Kempe connectivity is
reduced to a single local property:
\begin{definition}[Local Kempe reachability with fixed input, v23 Def.~8.12]
\label{def:lkr}
Let $X$ be a planar cubic tangle with ordered input stubs $S_{\mathrm{in}}$ and
output stubs $S_{\mathrm{out}}$, and fix an input coloring $c$. Let
$\Omega_X(c)$ be the set of proper 3-edge-colorings of $X$ extending $c$ on the
input stubs. The \emph{Kempe graph} $\mathsf{K}_X(c)$ has vertex set
$\Omega_X(c)$, with an edge whenever two colorings differ by a single Kempe
switch supported on a two-color component touching no input-stub edge. We say
$X$ satisfies $\LKR$ if $\mathsf{K}_X(c)$ is connected for every $c$ for which
$\Omega_X(c)\neq\emptyset$.
\end{definition}

The chain of Pillar~C is then:
\begin{itemize}
  \item \textbf{Lemma~8.14}: the canonical three-cell collar gadget $\tau$
        satisfies $\LKR$ (and so does its mirror);
  \item \textbf{Lemma~8.15}: two-color \emph{transparency} of $\tau$ --- every
        two-color component is a boundary-to-boundary path with no internal
        cycle;
  \item \textbf{Lemma~8.18 / Corollary~8.19}: $\LKR$ \emph{composes} across a
        transparent interface, so an entire cut-open radius-1 collar tangle
        $\widehat{A}=G_m\circ\cdots\circ G_1$ satisfies $\LKR$;
  \item \textbf{\S9.2 assembly}: peel the wide annulus $H$ into nested radius-1
        collars and use the collar $\LKR$ property to rewrite the current
        coloring to the target collar-by-collar from outside to inside, the
        cap-aware locality lemma (7.9) ensuring no switched component reaches
        back to change a fixed boundary.
\end{itemize}

\subsection{Why Lemma~8.14 is the keystone}

The manuscript states that its proof is human-readable and ``the only finite
verification required is a small by-hand audit of a single three-cell gadget''
(v23, Abstract and \S1). In the three-pillar architecture, every other step is
either an explicit local lemma or a structural reduction; the one place where a
finite, exhaustive check is genuinely unavoidable is the $\LKR$ property of the
canonical gadget $\tau$. Formalizing Lemma~8.14 therefore certifies exactly the
finite core that the manuscript isolates as load-bearing. The remainder of this
companion describes that formalization (\S\ref{sec:lkr}), the resolution of the
Pillar~A spanning question into a need-based repair (\S\ref{sec:pillarA}), the
Pillar~C composition ladder --- what is proved, what is reduced to one pin, and
where a decisive counterexample was found (\S\ref{sec:pillarC}) --- and a precise
status ledger (\S\ref{sec:status}). Superseded routes are recorded in
appendices.

\section{The Pillar-C keystone, machine-checked: Lemma~8.14 and transparency 8.15}
\label{sec:lkr}

We restate the manuscript's keystone in its own notation and then describe the
Lean encoding.

\begin{lemma}[$\LKR$ for the canonical three-cell gadget, v23 Lemma~8.14]
\label{lem:814}
Let $\tau$ be the canonical three-cell collar gadget. It is a planar cubic
tangle with internal vertices $F_0,\dots,F_5$ and boundary stubs
$B_0,\dots,B_7$, whose internal edges form the tree
\[
  F_1 - F_0 - F_2 - F_4 - F_5, \qquad \text{with a leaf } F_3 \text{ attached to } F_2,
\]
and whose boundary stubs are attached by
$F_0B_0,\,F_1B_1,\,F_1B_2,\,F_3B_3,\,F_3B_4,\,F_5B_5,\,F_5B_6,\,F_4B_7$. Take
input stubs $\{B_0,B_1,B_2,B_3\}$ and output stubs $\{B_4,B_5,B_6,B_7\}$. Then
$\tau$ satisfies $\LKR$. The same holds for the mirror gadget.
\end{lemma}

The manuscript proves Lemma~\ref{lem:814} by cutting $\tau$ along the internal
interface edge $I=F_2F_4$ into $\tau=\sigma_{\mathrm{R}}\circ\sigma_{\mathrm{L}}$
and doing a small case analysis on the two pieces. Our formalization instead
audits the \emph{full} gadget directly, by exhibiting the proper-coloring state
space and an explicit Kempe-path certificate connecting every fixed-input
fiber; this is a strictly more direct realization of the same $\LKR$ claim, and
it does not rely on the $\sigma_{\mathrm{L}},\sigma_{\mathrm{R}}$ split or on
the composition Lemma~8.18.

\subsection{The canonical gadget and its 192-state coloring space}

The Lean module fixes the colors and the gadget incidence exactly as in
Lemma~\ref{lem:814}: the three nonzero colors are an inductive type, and the
gadget's edges (five internal edges plus the eight boundary stubs) and internal
vertices are inductive types with an explicit incidence map.

\begin{lstlisting}[language={}]
inductive LColor | r | b | p
  deriving DecidableEq, BEq, Repr, Inhabited

inductive TauEdge
  | F1F0 | F0F2 | F2F3 | F2F4 | F4F5
  | B0 | B1 | B2 | B3 | B4 | B5 | B6 | B7
  deriving DecidableEq, BEq, Repr, Inhabited

def tauInputEdges : List TauEdge := [TauEdge.B0, TauEdge.B1, TauEdge.B2, TauEdge.B3]
\end{lstlisting}

A proper 3-edge-coloring of $\tau$ is represented by the finite choices left
after properness is imposed at every cubic vertex: an ordered color triple at
the central vertex $F_2$, one binary choice for the forced edge $F_1F_0$, and one
ordered-pair (swap) bit at each remaining boundary fork --- $6\times2^{5}=192$
states. That this enumeration is exactly the proper-coloring state space is a
checked fact ($192$ entries, no duplicates, proper at every vertex).

\subsection{The Kempe audit and the machine-checked theorems}

A Kempe step switches a two-color component \emph{disjoint from the input stubs}
(Definition~\ref{def:lkr}). The $\LKR$ property is witnessed by an explicit
indexed-path certificate: for each of the 192 states it records the
representative of that state's fixed-input fiber and a path of at most two
input-disjoint Kempe switches reaching it, plus a coherence audit that any two
states with the same input edges share a representative (so fibers are well
defined). The three finite audits each reduce to \texttt{true} by the kernel's
decision procedure; the only axiom dependency is \texttt{propext}.

\begin{lstlisting}[language={}]
theorem tauStateSpaceAudit_ok : tauStateSpaceAudit = true := by decide
theorem lemma814PathCertificate_rows_ok :
    validIndexedRowsFrom 0 lemma814PathCertificate = true := by decide
theorem representativeCoherenceAudit_ok : representativeCoherenceAudit = true := by decide

/-- Lemma 8.14, closed by the finite state-space and indexed path certificate. -/
theorem lemma814_tau_lkrIn_audit : lemma814_tau_lkrIn_finiteCheck = true := by
  simp [lemma814_tau_lkrIn_finiteCheck, tauStateSpaceAudit_ok, tauLKRInAudit,
    lemma814PathCertificateAudit_ok]
\end{lstlisting}

The mirror clause holds definitionally (reflection changes only the planar order
of an isomorphic incidence tree). The incidence enumeration is complete for the
encoded tau graph: every proper incidence coloring is one of the 192 enumerated
colorings. The finite certificate audit is proved, but graph-level soundness of
every encoded step remains the named obligation
\texttt{Lemma814EncodedMovesAreGraphKempeSwitchesTarget}. Thus the current
Lemma~8.14 result is a finite encoded-model audit, not yet the manuscript graph
theorem.

\begin{remark}[Transparency~8.15, finite encoded-model audit]
\label{rem:815}
The same module formalizes Lemma~8.15 (two-color transparency): for each of the
192 colorings, every nonempty two-color component is a boundary-to-boundary tree
path (\texttt{transparentComponent}), aggregated as
\texttt{tauTreeTransparencyAudit}. That audit is now \emph{proved} ---
\texttt{lemma815\_tau\_tree\_transparency\_audit} reduces to \texttt{true} by a
range decomposition of the finite check, axiom-clean. The encoded transparency audit is proved. Its use as graph-level composition
(\S\ref{sec:pillarC}) shares Lemma~8.14's named graph-fidelity obligation, so
the local graph-theoretic core is not yet sealed.
\end{remark}

\section{Pillar~A: a proper-subspace obstruction, and its need-based repair}
\label{sec:pillarA}

Pillar~A (the spanning Theorem~\ref{thm:49}) is where the formalization first
found --- and then routed around --- a real gap. We give both halves, because
together they are a clean example of what a companion of this kind is for.

\subsection{The obstruction: the spanning \emph{equality} over-reaches}

On every \emph{closed-walk exact shell} --- a concrete annular configuration
carrying the boundary-source and interior-endpoint witness data that
Theorem~\ref{thm:49} is meant to cover --- the boundary-zero target that the
Kempe-closure detector actually pins down is a \emph{strict} subspace of the full
boundary-zero chain space.

\begin{lstlisting}[language={}]
/-- On every closed-walk exact shell, the Theorem-4.9 boundary-zero target is a
proper subspace of the full selected-boundary-zero chain space. -/
theorem theorem49BoundaryZeroKirchhoffSubspace_lt_planarBoundaryZeroSubmodule_of_closedWalkExactShell
    [Fintype G.edgeSet] {emb : PlaneEmbeddingWithBoundary G}
    (shell : ClosedWalkExactShell emb) :
    theorem49BoundaryZeroKirchhoffSubspace emb
        (selectedBoundaryInteriorEdgeEndpointVertices emb) <
      planarBoundaryZeroSubmodule emb
\end{lstlisting}

The strict inequality \texttt{<} is the load-bearing content: a proved
\emph{proper}-subspace inclusion, not a heuristic. The structural reason (every
closed-walk shell forces two distinct interior-support edges onto one face
boundary) refutes the natural one-collar repairs, and a companion audit shows the
four proposed \emph{uniform geometric oracles} that would each close
Theorem~\ref{thm:49}'s synthesis are all false:
\begin{lstlisting}[language={}]
theorem not_interiorDualPeelOracle         : ¬ InteriorDualPeelOracle
theorem not_wellFoundedPeelOracle          : ¬ WellFoundedPeelOracle
theorem not_previousBoundaryGeometryOracle : ¬ PreviousBoundaryGeometryOracle
theorem not_residualBoundaryGeometryOracle : ¬ ResidualBoundaryGeometryOracle
\end{lstlisting}
Lean proves that the formal detector is strictly smaller than a larger
unrestricted selected-boundary-zero chain space. This does not refute v23's
actual cycle-space/Kirchhoff-constrained claim
$\Wz(H)\subseteq\mathrm{span}_{\FF}(\mathcal{G})$. That inclusion is now named
\texttt{V23Theorem49W0ContainedInKempeClosureGeneratorSpanTarget} and remains
open; the four uniform oracle routes that would have closed it uniformly are
refuted.

\subsection{The repair: the route never needs the full space}

The resolution is that the synthesis does not \emph{use} the full boundary-zero
space; it uses only the discrepancy classes of boundary-compatible coloring
pairs, and those lie in the detector subspace the manuscript's own peeling
witnesses span. Formally, the requirement actually consumed downstream is
\emph{need-based}: it suffices that (1) the detector subspace is contained in the
projected Kempe-closure generator span, (2) each execution step preserves that
subspace, and (3) every route-used discrepancy class lies in it. All three are
machine-checked in the route's natural cubic setting (module
\texttt{PillarASmallerSpace}):
\begin{lstlisting}[language={}]
-- (3) discrepancy classes lie in the detector subspace (cubic setting):
theorem routeUsedDiscrepancyClassesLieInDetectorSubspace_of_hasCubicIncidentEdgeTriples ...
-- (2) a Kempe-switch execution step preserves detector membership:
theorem routeKempeSwitchExecutionStep_preserves_detectorSubspace ...
-- assembled sufficiency (inclusion form), discharging the requirement:
theorem needBasedPillarARequirement_of_detector_le_span_of_stepPreserves_of_routeDiscrepancies ...
\end{lstlisting}
Proposition~(3) has a two-idea proof one can state aloud: a Tait coloring is a
Klein-four flow at each cubic vertex, so two colorings agreeing on the selected
boundary differ by a \emph{boundary-zero flow}, which is exactly what the
Kirchhoff detector certifies. The manuscript's Theorem~4.9 states an \emph{inclusion}, not an equality. The
existing need-based theorems are conditional on detector containment, step
preservation, discrepancy membership, and embedding-specific peel data. They
are conditional algebraic-span infrastructure, not a proof of Theorem~4.9. What remains a genuine per-embedding obligation is
the \emph{supply} of a peeling witness for the synthesis's actual embedding
class; that is a composition-stage task, not a spanning gap.

\section{Pillar~C: the composition ladder}
\label{sec:pillarC}

Between the local keystone (\S\ref{sec:lkr}) and 4CT, v23's own argument climbs a
ladder: single-gadget $\LKR$ $\to$ transparent composition to a whole collar
chain (Gate-2) $\to$ closed re-glued collars (\S9.2 Step~3) $\to$ execution in the
ambient annulus (\S9.2 Step~4) $\to$ the wide-annulus peel. This companion has
now walked a substantial part of that ladder, with one step proved-modulo-a-pin
and one step \emph{refuted}.

\subsection{Gate-2: transparent composition, reduced to one pin}

Corollary~8.19 lifts $\LKR$ from the single gadget to a whole cut-open radius-1
collar chain across a transparent interface. The manuscript's own Step-1 proof of
Lemma~8.18 (a pointwise preparation via Lemma~8.16) does \emph{not} go through ---
the formalization refutes it directly:
\begin{lstlisting}[language={}]
theorem lemma818_pointwise_lift_preparation_obstruction : ...   -- axiom-clean
\end{lstlisting}
In its place the development builds a frontier-automaton / fibration architecture:
a finite frontier state over the chain transcript, generated parent-row
certificates, and a fibered-connectivity decomposition (base connected + fibers
connected + lifts). The executable chain test \texttt{chainLKRInAudit} passes on
all words up to length seven (real-fiber labs, all input fibers connected, zero
failures). The former metadata-only classification audit has been replaced: the
actual per-target test is
\texttt{representativeTargetFiberKempeOrbitConnectivityAudit}, and Gate-2's
single live pin is
\begin{lstlisting}[language={}]
Gate2RepresentativeFiberKempeOrbitConnectivityPin
\end{lstlisting}
combining fourteen representative real-fiber audits with DFA-mode sufficiency;
the all-chain implication is proved conditional on this pin. (A separate refuted route,
\texttt{not\_concreteChainFiberAppendPrefixFibrationClosed}, records a shortcut
that provably fails and is kept dead.)

\subsection{S2: length-2/3 radial-closing-edge fiber census}

Gate-2 certifies $\LKR$ with the \emph{first} gadget's input stubs fixed;
\S9.2 Step~3 needs the \emph{cyclically re-glued} collar with all outer-side stubs
fixed --- different fibers, no bridging lemma in the manuscript. What is proved
is a finite census: it ranges over twelve tau/mirror words of length 2 or 3 and
fixes the four radial-closing edges (a lab found all 972 such fibers
connected-or-empty --- 216 nonempty, 756 empty, zero disconnected --- and the
census target was discharged with zero pins). It does \emph{not} fix every
outer stub and does not range over arbitrary collar length. The actual
manuscript claim is named
\texttt{ManuscriptSection92S2AllOuterStubsFixedArbitraryCollarTarget} and
remains open.
\begin{lstlisting}[language={}]
def closedCollarLength23RadialClosingEdgeFiberCensus : Prop := ...   -- proved
def ManuscriptSection92S2AllOuterStubsFixedArbitraryCollarTarget : Prop  -- open
\end{lstlisting}
The route survey had rated S2 the sharpest mismatch in the manuscript --- a
statement never written down in v23 and not implied by Gate-2 --- so proving it on
the finite class is a real advance.

\subsection{S4: abstract-chain ambient-escape regression}

\S9.2 Step~4 reinterprets a collar Kempe word as a word of the \emph{ambient}
annulus, and confines the switched components with a face-only locality lemma
(7.9). The formalization proves an ambient-escape regression in the encoded
tau-chain model --- a collar plus one deeper radial ring sharing the interface.
No theorem currently realizes that witness as a simple planar,
cyclically-five-edge-connected actual collar; graph-facing salvage is
conditional on
\texttt{ClosedWalkActualCollarRotationFaceIncidenceCutLiftObligation}, which
remains open.
\begin{lstlisting}[language={}]
/-- A collar-certified, collar-input-avoiding Kempe switch whose ambient
two-layer component leaves through the inner side, hits the fixed outer input,
and does NOT restrict to the certified collar move. -/
theorem section92Step4TwoLayerCounterexampleClass :
    Section92Step4TwoLayerCounterexampleClass         -- axioms [propext]
\end{lstlisting}
Concretely (outer word $\tau\tau$, inner word $\tau$, input key $r,r,b,r$): a
switch that is input-avoiding \emph{in the collar alone} balloons, in the
two-layer object, into a component that crosses the shared inner boundary into the
deeper ring, returns, meets the fixed outer input edge, and lands the collar on a
\emph{different} coloring than the certified move. This refutes the reinterpretation
as a regime-free, purely-local principle: the ambient switch is a \emph{different}
move on the collar than the one Gate-2/S2 certified.

We are precise about what this does and does not say. It is a defect in the
manuscript's assembly \emph{seam} as argued (Step~4 invokes 7.9, a lemma proved
for single-face collars, as if it were a global confinement), not a refutation of
4CT.

The cheapest repair has also been tested, and \emph{refuted}. Strengthening the
certification so that switches avoid the entire outer interface --- both input
\emph{and} output stubs, not only the inputs --- does eliminate the escape (a
two-layer companion run finds zero escapes under full-interface avoidance), but at
a fatal cost: on the same finite class, two colorings with the same fixed input
are connected \emph{only} by a switch whose component touches the output stubs, so
forbidding output touches disconnects the fibers (module
\texttt{GoertzelLemma818TwoSidedAvoidance}: $432$ open and $216$ closed nonempty
fibers disconnect). The escape channel is therefore \emph{load-bearing for
connectivity}: the very moves that let a component escape the collar are the moves
the fibers need to remain connected. Confinement and connectivity are in direct
tension, so any repair must confine the escaping components \emph{without}
severing collar connectivity --- a needle the simplest fix cannot thread.

What remains is a genuinely regime-level salvage: a proof that the
minimal-counterexample \emph{regime} (CAP5-free normal form) excludes the
disconnecting/escaping configurations, or a \S9 redesign that only ever peels
\emph{face} collars where 7.9 applies as stated. Neither is established; both need
planar-normal-form infrastructure the development does not yet have; and the
abstract gadget lab cannot decide them, since the tension may or may not dissolve
in the actual normal form. As of this writing the seam is open, with both the
escape and the failure of the cheapest repair pinned down as checked theorems.

\subsection{Sharpening the open salvage: a winding invariant and a cyclic 2-cut reduction}

Since the escape was pinned down, the formalization has \emph{sharpened} --- though
not closed --- the open salvage, by asking a Spencer--Brown/Kauffman-native
question: which two-color components can escape? On the closed annular collar a
two-color component carries a well-defined mod-2 \emph{winding class} around the
annulus generator --- the parity of its crossing count with a fixed radial cut ---
a Parity-Lemma-style invariant, well defined per component and tracked across Kempe
switches. In this language the escape is exactly a \emph{winding} move: the
components that leave through the inner side and return are those of nonzero winding
class. The precise form of the S4 freedom is then a homological statement,
machine-checked on the witness class: \emph{winding is not boundary-determined} --- a
fixed-boundary fiber can carry components of different total winding class, so the
escaping (winding-1) moves are not forced by the boundary data alone.

This localizes the difficulty to one geometric question: can a winding-freedom
escape be \emph{realized on a genuinely simple planar collar}? The formalization
reduces it to a classical connectivity fact. Any simple planar realization that
preserves the winding profile must replace the escape's parallel endpoint bundle by
a simple two-pole side; a tree-rank (handshaking) count ---
$2\,|\text{side edges}| = 3\,|\text{internal}| + 4$ forces
$|\text{internal}|+1 < |\text{side edges}|$, hence a cycle --- turns that side into
a \emph{cyclic 2-cut} of the collar. A collar admitting no realized cyclic 2-cut
therefore admits no escape realization:
\begin{lstlisting}[language={}]
-- every simple planar winding-escape realization forces a cyclic two-cut:
theorem ...Escape_forces_cyclicTwoCut_allSizes : ...
-- so a collar forbidding cyclic two-cuts (e.g. cyclically 5-edge-connected)
-- admits no escape realization:
theorem ...Escape_not_simplyRealizable_allSizes
    (h : ...ForbidsCyclicTwoCut G) : ¬ ...SimplePlanarEscapeRealization G
\end{lstlisting}

We are precise about what this buys. It is a \emph{reduction of the open salvage to
one named residual}, not a repair: the theorem is conditional on the collar
forbidding cyclic 2-cuts, and the current realization datum carries the
separating-pair structure as an obligation rather than deriving it from the collar's
embedding. What it does is convert the vague ``the minimal-counterexample regime
should exclude the escape'' into a concrete checkable target --- \emph{the escape
collar carries no cyclic 2-cut}, i.e.\ it is cyclically sufficiently edge-connected,
which is exactly the kind of property the minimal-counterexample normal form is
expected to supply. And this residual is not an un-formalizable abstraction: the
development already proves cyclic-five-edge-connectivity for concrete graphs (and
refutes it for others), so discharging or refuting it \emph{for the escape collar}
is a reachable next question --- the honest current frontier of the S4 salvage.

\section{The v24 route: deserts, pivots, and the necklace}
\label{sec:v24}

Goertzel's v24 working package (July 2026) replaces v23's \S6--7 with a new
architecture: pentagonal \emph{pivots}, hexagonal \emph{deserts}, a
\emph{necklace} decomposition of the sphere relative to the twelve pentagons,
winding and monodromy invariants along corridors, and a final \emph{Trail
Completability} target from which the minimal-counterexample contradiction is
to follow, conditional on a dependency chain L1--L10. The package proposes a
universal Short Strand Conjecture and then refutes it itself with a
$688$-edge, reach-$40$ strand, so the route requires a replacement bound.
This section records what the formalization has established of this route so
far; the formalization targets combinatorial rotation systems throughout, not
topological planarity.

\paragraph{Independent reproduction of the validation toolchain.} Every
numerical claim consumed from the package was recomputed independently before
use: the annular censuses ($5412$ and $31704$) with their Kempe fibers, the
fullerene finite bases, the pentagon finite seals, tube seed growth, the
desert and strand specimens (including the Short-Strand refuter), Goldberg
normal-form specimens, and the GWCO exact base through twelve. A recurrent
tube-ring transfer theorem replaces the package's Pascal-triangle
extrapolation (\texttt{GoertzelV24TubeRingTransfer}).

\paragraph{Rotation systems and facial linear algebra.} A Tait-coloring
layer is built over rotation systems, with each edge's two faces identified,
face simplicity reduced to two-sidedness, facial dual connectedness, and the
exact facial linear algebra: facial boundary rank, primal cycle-space rank,
the equality of the facial span with the cycle space, and the face winding
classification (\texttt{GoertzelV24FaceBoundaryLinearAlgebra},
\texttt{GoertzelV24OrbitFaceCycleSpaceEquality},
\texttt{GoertzelV24WindingClassification}).

\paragraph{Curvature.} The all-face discrete Gauss--Bonnet identity is
proved from rotation-system data alone: the total face curvature
$\sum_f\,(6-\ell(f))$ equals $12$ (\texttt{totalFaceCurvature\_eq\_twelve} in
\texttt{GoertzelV24CurvatureScope}). The package's \emph{at-most-twelve
pentagon} step, which is false as stated under minimum face size five, is
corrected by a proved \emph{weighted} pentagon budget charging faces by
$6-\ell(f)$; the fullerene hexagonal bulk is derived on orbit faces, and
curvature-to-bulk corridor bounds connect the budget to corridor extraction
(\texttt{GoertzelV24OrbitFaceCurvatureBulk}, \texttt{GoertzelV24BulkCorridor}).

\paragraph{Necklace and corridors.} The locked-pentagon necklace is
assembled and the singleton necklace strand forced
(\texttt{GoertzelV24Necklace}). A hexagonal corridor theory is developed:
defect-free corridor neighborhoods extracted from geodesic face paths; an
internal rung classification with induced rung chains and the exclusion of
adjacent induced rungs; clean width-four slabs with oriented interfaces and
the matching of consecutive interfaces; Tait transfer across slabs with
finite color-transition alphabets and color-state words; and a sharp
chirality barrier under which transfer across a chiral reversal forces the
melted state (\texttt{GoertzelV24CleanHexCorridor},
\texttt{GoertzelV24HexFaceRungType}, \texttt{GoertzelV24OrientedHexSlab},
\texttt{GoertzelV24HexCorridorColorTransfer},
\texttt{GoertzelV24HexCorridorChiralityBarrier}).

\paragraph{Pumping.} The corridor acceptance language is proved periodic and
corridor pumping congruence-correct; with the recurrent tube-ring transfer
these are the working tools of the Short-Strand replacement search
(\texttt{GoertzelV24CorridorPumping}).

\paragraph{Framed trails and local fiber calculi.} The framed-trail layer
proves completion semantics and defect freedom --- the unique admissible
missing color and the characterization of proper trail extensions
(\texttt{missingColorAt\_is\_unique\_admissible},
\texttt{hasMatchingDefectColors\_iff\_exists\_properTrailExtensionColor} in
\texttt{GoertzelV24FramedTrail}) --- supported by exact local calculi: the
square fiber table with its realized Kempe switch and local and
compositional profile identities, the digon Kempe rung, rotor local color
routes, and the two-edge-cut color parity (\texttt{GoertzelV24Square}\,%
/\texttt{SquareKempe}/\texttt{SquareComposition}, \texttt{GoertzelV24Digon},
\texttt{GoertzelV24Rotor}, \texttt{GoertzelV24TwoEdgeCut}).

\paragraph{The minimality normal form (proved).} The full minimality
connectivity ladder is now machine-checked on rotation systems, with no
\texttt{sorry} and axioms in
$\{\texttt{propext},\texttt{Classical.choice},\texttt{Quot.sound}\}$. The
two-edge, cyclic three-edge, and cyclic four-edge cut reductions are each
proved by capping both sides into strictly smaller cubic rotation systems and
gluing their Tait colorings across the cut; a vertex-minimal Tait
counterexample admits none of them, so it is cyclically five-edge-connected
(\texttt{cyclicallyFiveEdgeConnected\_of\_vertexMinimalTaitCounterexample}).
The four-edge rung is the classically hard one: at a four-edge cut the four
boundary colors sum to zero in the Klein four-group, giving three boundary
classes (all-equal and the two non-crossing adjacent pairings), and the two
capped sides always share a reducible class. The load-bearing step --- that a
side admitting an all-equal boundary escapes, by a Kempe swap, into a
non-crossing adjacent class (never the crossing Birkhoff pairing) --- is
\emph{proved}, not assumed, from the planarity of the interior Kempe chains
(\texttt{GoertzelV24FourEdgeCutKempeEscape}), and a minimal counterexample is
correspondingly forced to have minimum face size five. The corridor-existence
side is lifted off the fullerene premise entirely: on any spherical cubic map
of minimum face size five, the pentagon count equals $12$ plus the total
negative curvature weight, and either the face count lies below a
weight-dependent threshold or a clean hexagonal corridor exists
(\texttt{GoertzelV24WeightedOrbitFaceCorridor}) --- no restriction to
pentagonal and hexagonal faces is used anywhere.

\paragraph{In progress and open.} On this normal form the endgame reduction is
under active construction: a fusion-chain / rebase-pumping dynamics that, from
any adjacent pair of the counterexample, produces either a fusion terminal or
a strictly descending rebase pumping
(\texttt{exists\_rotationOrderedFusionTerminalOrRebasePumping\_of\_}%
\texttt{vertexMinimalTaitCounterexample}), together with an exhaustive scope
classification of the remote facial-dual cycles it generates (length $\geq5$
exact cut, length-three single-vertex cut, length-four adjacent-pair cut with
an internal dual chord). What remains \emph{open} is the coloring transport
that turns each classified case into a contradiction with minimality --- the
large-cut splice/interface law and the explicit consumption of the
singleton-star and adjacent-pair-collar cases --- and hence Trail
Completability, the L1--L10 chain, and a proved Short-Strand replacement. The
v23-era collar obligations are taken up only where this route consumes them.
Nothing in this section claims the Four-Color Theorem or the completion of the
v24 route.

\paragraph{Canonical data at reentry joins.}
The rotation ordering removes a small but potentially serious compositional
ambiguity.  If two adjacent-pair presentations have the same ordered central
edge, then their rotation-ordered four-port data are equal, regardless of the
original enumeration of the four boundary ports.  We therefore refine the
recursive relation so that each step retains its source fusion normal and its
certified pumping launch.  Given two consecutive steps, the target Kempe
package of the first supplies a normal with a genuine same-base Kempe escape,
and the second supplies its source normal; the two recovered rotation-ordered
data agree exactly.  Finite iteration consequently yields a structural
terminal, a large exact remote separator, or a nonempty closed orbit whose
successive normal forms are canonically aligned.  This alignment is a
bookkeeping theorem, not yet a coloring-transport theorem: transport between
different deleted adjacent pairs remains the next local collar problem.

\section{Status and scope: machine-checked, reduced, and open}
\label{sec:status}

\paragraph{Historical snapshot.}
This section records the companion's 21 July 2026 v23/v24 status.  It is kept
for provenance and should not be read as the present closure ledger; the
controlling August status is Section~\ref{sec:current}.

We state the scope exactly. The full route is a ladder $S0,\dots,S9$ (Gate-2, then
v23 \S9 and \S3); the table records each step's manuscript reference and current
Lean status.

\begin{center}\small
\begin{tabular}{@{}llll@{}}
\toprule
Step & Content & v23 ref & Lean status \\
\midrule
--- & single-gadget $\LKR$                & Lem.~8.14        & finite encoded-model audit; graph move-soundness open \\
--- & two-color transparency              & Lem.~8.15        & finite encoded-model audit; shares graph-fidelity boundary \\
--- & Pillar~A                            & Thm.~4.9         & actual $\Wz$ inclusion open; unrestricted-oracle strictness proved \\
S0  & Gate-2 chain $\LKR$                  & Lem.~8.18/8.19   & one actual fiber/orbit plus mode-sufficiency pin \\
S2  & radial-closing-edge fiber census     & \S9.2 Step~3     & proved only at length 2/3; manuscript target open \\
S4  & ambient escape regression            & \S9.2 Step~4     & abstract-chain countermodel; graph realization open \\
S1  & collar $\Rightarrow\tau$/mirror chain & Lem.~8.5--8.6   & open (structural) \\
S3  & annulus peeling well-definedness     & \S9.2 Step~2     & open (structural) \\
S5  & interface to Kauffman Thm.~3.1       & \S9.2 / \S3      & open (statement-level) \\
S6  & Kauffman parity/primality internals  & \S3.2--3.3       & open (structural) \\
S7  & minimal-counterexample normal form   & Lem.~5.2--5.3    & open (structural) \\
S8  & CAP5 reducibility                    & \S6              & partial \\
S9  & minimality wrap ($+$ Tait)           & Def.~5.1 / Thm.~2.3 & open (structural) \\
\bottomrule
\end{tabular}
\end{center}

\paragraph{Machine-checked outright.} The finite encoded-model audits for
Lemmas~8.14 and 8.15, together with completeness of the 192-state incidence
enumeration; unrestricted-oracle strictness and conditional need-based Pillar~A
infrastructure (with four uniform peel oracles refuted); and the length-2/3
radial-closing-edge S2 census. The corresponding graph move-soundness, actual
v23 $\Wz(H)$ inclusion, and manuscript S2 statement remain open. All with no \texttt{sorry} and axioms
$\subseteq\{\texttt{propext},\texttt{Classical.choice},\texttt{Quot.sound}\}$.

\paragraph{Reduced to one pin.} Gate-2 (Lemmas~8.18/8.19): the manuscript's
pointwise Step-1 preparation is refuted and replaced by the frontier-automaton /
fibration architecture, whose real-fiber labs pass to length seven; the whole
composition rests on a single named structural remainder.

\paragraph{Regression countermodel.} \S9.2 Step~4's collar-word-to-ambient-word
reinterpretation fails in the encoded two-layer tau-chain model. A simple planar
cyclically-five-connected graph realization has not been proved; the
actual-collar cut-lift obligation remains open. In the encoded model this
lowers the odds that the v23 route completes \emph{as written}; its salvage is
open.

\paragraph{v24 scope.} See \S\ref{sec:v24}. The validation toolchain and
censuses are independently reproduced; the all-face curvature identity and
the corrected \emph{weighted} pentagon budget are proved; the
necklace/corridor, pumping, and framed-trail layers are in the tree; the
endgame (fusion chains, chiral monodromy, Trail Completability, L1--L10, and
a proved Short-Strand replacement) is open.

\paragraph{Open (the non-finite heart).} Steps S1, S3, S5--S9: the realization of
abstract chains as actual collars, the wide-annulus peeling, the interface to
Kauffman's reduction, and the Kauffman/normal-form/CAP5 internals. Their common
cost is one shared planar-embedding/duality/Jordan infrastructure the development
does not yet have; most carry a known proof shape, so their risk is
formalization cost rather than mathematical doubt --- with the exception of the S4
salvage, which is genuinely open.

\paragraph{What is not claimed.} We do \emph{not} claim the Four-Color Theorem,
and we do not claim that v23 proves it. The 4CT is a theorem via the
computer-assisted proofs; what is at issue is solely the viability of the v23
human-readable route, which remains \emph{open} --- its foundational layer now
machine-checked, its most dangerous assembly seam pinned down as broken.

\section*{Provenance}
\addcontentsline{toc}{section}{Provenance}
The August revision preserves the accumulated July audit and its negative
results, and adds Section~\ref{sec:current} as the controlling status for the
current Goertzel/Fable reductive--compositional development.  Future revisions
should update that section when a closure obligation changes status and append
failed constructions to Appendix~\ref{app:graveyard}, rather than erasing the
reason they were abandoned.

The manuscript discussed is Ben Goertzel, \emph{A Human-Readable,
Spencer--Brown/Kauffman-Style Proof of the Four-Color Theorem}, version~v23. The
formalization is work of the MeTTapedia formalization effort. The keystone and
transparency are module \texttt{GoertzelLemma814}; the Pillar-A obstruction and
its need-based repair are \texttt{FrontierGeometric}, \texttt{GoalAudit}, and
\texttt{PillarASmallerSpace}; the Gate-2 architecture is
\texttt{GoertzelLemma818SemanticBridge} and the frontier-automaton modules; the
S2 bridge is \texttt{GoertzelLemma818ClosedCollarS2Bridge}; the S4 counterexample
and the refuted two-sided-avoidance repair are
\texttt{GoertzelLemma818TwoLayerEscape} and
\texttt{GoertzelLemma818TwoSidedAvoidance}. The v24 route of \S\ref{sec:v24}
is the \texttt{GoertzelV24*} module family (roughly sixty modules; the primary
mathematical source is the v24 working package \emph{Deserts, Pivots, and the
Necklace}, July 2026). Lean listings are ASCII
transliterations of the source declarations (the sources are normative); every
headline was checked with \texttt{\#print axioms}. Per-declaration source links
are omitted pending publication of the development tree.

\appendix

\section{A graveyard of superseded sub-routes}
\label{app:graveyard}

Each entry is a route or reading that was tried and set aside; several are
themselves machine-checked theorems that a stated step \emph{fails}.

\begin{description}
  \item[The spanning \emph{equality} $\Wz(H)=\mathrm{span}(\mathcal{G})$.] Refuted
        on every closed-walk exact shell (\S\ref{sec:pillarA}); superseded by the
        need-based inclusion, which is what the synthesis actually consumes.
  \item[The four uniform geometric peel oracles.] Each would have closed
        Theorem~4.9's synthesis uniformly; all four are machine-checked false.
        The route uses the manuscript's own per-embedding peeling witnesses
        instead.
  \item[The pointwise Step-1 preparation of Lemma~8.18.] The manuscript's own
        pointwise lift (via Lemma~8.16) is refuted
        (\texttt{lemma818\_pointwise\_lift\_preparation\_obstruction}); replaced by
        the frontier-automaton / fibration composition.
  \item[The prefix-fibration shortcut for Gate-2.] A tempting closed-form
        shortcut, machine-checked to fail
        (\texttt{not\_concreteChainFiberAppendPrefixFibrationClosed}); kept dead.
  \item[The false ``additivity'' of Kempe switches.] An $\FF$-additive execution
        of the Pillar-A decomposition $\Delta=\sum v_{f,g}$. This is \emph{not}
        the manuscript's route --- Pillar~C is titled ``executability without
        false additivity,'' and v23 \S8.1--8.3 show mixed color-pair switches on
        shared support do not commute. Recorded only to set aside; the faithful
        route is the operator/$\LKR$ execution of \S\ref{sec:lkr}.
  \item[An $\FF$-additivity / evader reconstruction (disavowed).] An early
        exploratory reconstruction outside the manuscript's route, which v23
        explicitly disavows (additivity invalid for mixed color-pairs; Rmk.~8.10).
        Not part of this companion's formalized content.
  \item[Boundary color-word agreement as full profile agreement.] Refuted by a
        two-state off-diagonal witness in
        \texttt{GoertzelV24ProfileDiagonalLiftFalsifier}: matching projected
        color support need not supply a diagonal entry for the full profile.
        The source splice therefore uses the finite full-profile diagonal
        audit, not its color projection.
  \item[Consecutive six-edge slab boundaries as radial splice cuts.] Refuted by
        \texttt{sourceTwoTileAlignedBoundaryWalks\_consecutive\_not\_disjoint}.
        Consecutive local boundary walks share their intervening dual
        interface.  This rules them out as the two separated annular cuts while
        retaining them as useful local windows.
  \item[A single generic hole crosscut as the corridor splice.] Generic dual
        connectedness produces one simple hole-to-hole slit.  It supplies
        neither a second disjoint slit nor alignment with the selected Cell--3
        corridor, so promoting it to the source splice would change the
        statement rather than prove it.
\end{description}

\end{document}
